\begin{abstract}
Consumer-progress acknowledgments in Redis Streams and NATS JetStream have different persistence implications even under always-synchronize append policies. We evaluate that distinction alongside recovery timing, consumer concurrency and publication cost using single-replica deployments on a two-node Kubernetes homelab. The study compares Redis 7.4.2/8.10.2 and NATS Server 2.10.24/2.15.0, preserving an original campaign and adding three deployments with fresh broker processes and stores, matched publication durations, CPU sensitivity and storage diagnostics. All 240 active follow-up recovery episodes completed, including live, deliberately unacknowledging consumers; twelve plain Redis read controls remained pending. Receipt timing followed residual eligibility plus policy-dependent excess. Redis's integrated CLAIM path exhibited greater excess than the tested explicit reclamation loop in this quiet, one-message workload. Sixteen consumers increased current memory/periodic drain throughput by 10.43--14.49 times across deployment-specific comparisons, while cycle p99 increased and driver capacity affected rates. At sixteen publishers, current periodic/memory publication ratios averaged 0.8257 for Redis and 0.8207 for JetStream. Ten of 72 planned always-profile publication trials failed (13.9\%), including four warm-up failures, compared with one of 72 periodic trials and none of 72 memory trials. Timed failures returned five-second client timeouts with unknown confirmation outcomes. A compact analytical model and selected Lean-checked statements delimit the claims; full proofs and detailed results are included in the appendices. The licensed artifact retains timings, failures, configurations and reproducible analysis. These observations characterize specified policies and acknowledgment costs; they do not establish equal durability contracts, power-loss survival or hardware-independent rankings.
\end{abstract}
